% ================================================================
%  chapters/parts/ch08-b.tex   （第 8 章 核磁共振方法，分片 b）
%  源：S2（下册）第 8 章，书内 p.540—550 / PDF p.32—42
%  覆盖：8.3 固体 NMR 在催化剂酸性研究中的应用（8.3.1—8.3.3）
%        8.4 催化剂表面吸附分子的 NMR 研究（8.4.1 孔道内有机物；8.4.2 129Xe）
%  未覆盖（分片 c）：8.5 原位 MAS NMR、8.6 积炭（8.4.2 第 5 小节的积炭细节亦在彼处）
%  页码说明：下册印刷页 526 未出现在 PDF 中，书内页码跳号属正常，一律以分片源文
%            的 [[DOWN PDF p.n | 书内 p.m]] 标记为准。
%  注：本分片源文 OCR 层\textbf{系统性丢失负号}（如 −68 记作 68）。凡负号均按上下文
%      （物理吸附 TMP 应在高场、L 酸络合物在 −60～−32 区间）恢复，并在相应位置用
%      \tocheck 标明"符号系恢复"，供集中核验。
% ================================================================

\subsection{固体 NMR 在催化剂酸性研究中的应用}
\label{sec:acidity}

\begin{probout}
本节要输出四个彼此独立的量：\textbf{酸的类型}（B 酸还是 L 酸）、\textbf{酸强度}、
\textbf{酸密度}（中心数目）与\textbf{可接近性}（酸性位在内表面还是外表面）。与常用方法相比，
红外光谱受消光系数制约，把吸收峰强度换算成基团数目存在系统性困难，峰重叠又使分辨率下降；
\ce{NH3} 程序升温脱附（\ce{NH3}-TPD）测的是脱附行为，\textbf{分辨不了酸的类型}。固体 NMR 有
两条入口：一是用 $^{1}$H MAS NMR 直接看桥羟基（B 酸位）；二是先让探针分子
（\ce{(CH3)3P}、\ce{^{15}N}-吡啶、\ce{^{13}CO} 等）吸附，再测探针核的化学位移。后者不仅
能区分酸的类型、强度与密度，借助\textbf{不同尺寸}的探针分子还能获得酸性位的空间信息。
\src{S2 §8.3，书内 p.540 / PDF p.32}
\end{probout}

\begin{principle}
\textbf{用碱性探针测酸强度的统一逻辑}是一条三步因果链：

\begin{center}
\fontsize{9}{13}\selectfont
碱的孤对电子（或 $\pi$ 体系）与酸位成键／成氢键 $\rightarrow$ 被观测核周围电子云被拉走
（去屏蔽）$\rightarrow$ 化学位移向\textbf{低场}（$\delta$ 增大）移动，\textbf{位移量随相互作用
强度单调增加}
\end{center}

由此得到三条使用纪律。（i）\textbf{位移是相对标尺，不是绝对酸强度}：只有\emph{同一探针、
同一参考态、同一温度}下的位移才能相互比较，标尺必须用一组酸强度已知的固体（SAPO-5、
HY、HZSM-5、HZSM-22 直到 \ce{AlCl3}/100\% \ce{H2SO4}、\ce{SbF5}）标定。（ii）\textbf{位移只
报告探针访问得到的酸位}：进不去孔道的探针不产生信号——这既是漏检的风险，也是可以选择
探针尺寸、只测外表面的手段。（iii）\textbf{位移给强度，峰面积给数目}，两者的成立条件不同，
后者的门槛见 8.3.3 的定量纪律。
\end{principle}

\subsubsection{$^{1}$H MAS NMR：不用探针的直接法}
\label{sec:h1-acidity}

分子筛的 Brönsted 酸位与骨架桥式羟基（Si—OH—Al）相关联，因此 $^{1}$H MAS NMR 是研究
表面 B 酸位\textbf{最直接}的方法——不需要引入任何探针分子，也就不存在"探针进不去孔道"
的问题。代价是固体中质子偶极--偶极作用很强、质子的化学位移范围又小，谱线极宽。可用的
窄化手段有：高速魔角旋转、特殊脉冲序列（如 \term{CRAMPS}{combined rotational and
multiple-pulse spectroscopy}{联合旋转与多脉冲谱}）、用 $^{2}$H 同位素稀释表面质子，以及
对催化剂表面做充分预处理。\src{S2 §8.3.1，书内 p.540 / PDF p.32}

\begin{measure}
\textbf{装样方式决定了能跑多快，也决定了谱可不可信。}酸性研究前样品必须\textbf{真空脱水}，
再装入转子或吸附探针。早期用玻璃封管法：催化剂在真空系统中抽空、吸附气体后用火焰密封，
制成小安瓿；玻璃管很难做得均匀对称，装入转子后转速一般低于 \qty{4}{kHz}。Haw 等设计的
\term{CAVERN}{}{原位处理装样装置}把催化剂平铺在床层上，使加热脱水和吸附都均匀进行，装样
后通过连杆密封转子，能有效防止装样时引入水，并允许较高转速；类似功能的装置用捣杆原位
装样、支架帽密封，转子最高转速可达 \qty{12}{kHz}。\src{S2 §8.3.1，书内 p.540 / PDF p.32}
\end{measure}

分子筛表面 OH 基的 $^{1}$H 化学位移归属如下（\src{S2 §8.3.1，书内 p.541 / PDF p.33}）：

\begin{itemize}[leftmargin=1.6em,itemsep=1pt,topsep=2pt]
\item $0.5\sim-0.5$ ppm：分子筛外表面或游离于骨架外、与金属离子相连的羟基（M—OH）；
\item $1.2\sim2.2$ ppm：缺陷位或端基硅羟基（Si—OH）；
\item $2.8\sim3.6$ ppm：与骨架铝相连、并与周围氧形成氢键的非骨架铝羟基（Al—OH）；
\item $3.6\sim4.3$ ppm：位于大笼或孔道中的桥式羟基（Si—OH—Al），即主 B 酸位；
\item $4.6\sim5.2$ ppm：位于八面沸石型分子筛小笼中的桥式羟基；
\item $5.2\sim7.0$ ppm：HZSM-5 与 H$\beta$ 中与骨架有静电作用的\textbf{第二种}桥式羟基。
\end{itemize}

\begin{principle}
\textbf{位移与红外波数之间的经验关系。}Brunner 把 $^{1}$H 化学位移 $\delta_{\mathrm{H}}$ 与该
OH 基在红外谱中的伸缩振动波数 $\nu_{\mathrm{OH}}$ 关联起来，给出两个经验式
（\src{S2 §8.3.1，书内 p.541 / PDF p.33}）：

\begin{equation}
\label{eq:brunner-iso}
\delta_{\mathrm{H}}\ (\mathrm{ppm})=57.1-0.0147\,\nu_{\mathrm{OH}}\ (\mathrm{cm^{-1}})
\qquad \text{（孤立 OH）}
\end{equation}

\begin{equation}
\label{eq:brunner-hb}
\delta_{\mathrm{H}}\ (\mathrm{ppm})=37.9-0.0092\,\nu_{\mathrm{OH}}\ (\mathrm{cm^{-1}})
\qquad \text{（存在氢键的 OH）}
\end{equation}

两式都写成"截距 $-$ 斜率 $\times\nu_{\mathrm{OH}}$"：波数越低（氢键越强、O—H 越弱），
$\delta_{\mathrm{H}}$ 越大，即向低场移动。\emph{核对说明}：已对照原页面图像
（书内 p.541 / PDF p.33），原书式（8-3）（8-4）逐字印作
$\delta_{\mathrm{H}}\,(\mathrm{ppm})=57.1-0.0147\,\nu_{\mathrm{OH}}\,(\mathrm{cm^{-1}})$、
$\delta_{\mathrm{H}}\,(\mathrm{ppm})=37.9-0.0092\,\nu_{\mathrm{OH}}\,(\mathrm{cm^{-1}})$，
减号、斜率与有效位数均与此处一致（源 OCR 文本层丢失了减号）。原书同时指出：
对与 Mg$^{2+}$、Ca$^{2+}$ 等金属离子相连的 OH，两式有偏差。
\end{principle}

\pedset{教学补充的自洽性核算（非原书数据）：把 HZSM-5 桥羟基的红外波数
$\nu_{\mathrm{OH}}\approx\qty{3610}{cm^{-1}}$ 代入式~\eqref{eq:brunner-iso}，得
$\delta_{\mathrm{H}}\approx4.0$ ppm，与上表 $3.6\sim4.3$ ppm 的桥羟基归属自洽；把原书提到
的 $\qty{3250}{cm^{-1}}$ 宽带代入式~\eqref{eq:brunner-hb}，得 $\delta_{\mathrm{H}}\approx8.0$ ppm，
与原书记录的 7.0 ppm 宽峰相差约 1 ppm——原书称"两个值能很好地符合"，读者可用同一组
公式复核这一对的偏差量。}

\begin{dataflow}
\textbf{从"看不见的铝"到脱羟基机制。}Freude 等对 HNaY 脱羟基过程做 $^{1}$H 与 $^{27}$Al
MAS NMR 定量，结果是桥羟基浓度的降低量与四配位骨架铝的减少量\textbf{相等}；这与
Uytterhoeven 等提出的"两个桥羟基缩合脱水生成一个骨架三配位铝"的模型不符。邓风等在不同
温度焙烧的 HZSM-5 上进一步发现两种脱羟基机制并存：一部分脱羟基伴随脱铝，另一部分只脱
羟基而不脱铝——后者的铝仍留在骨架上，但四极作用增强而变得"不可观测"。
\src{S2 §8.3.1，书内 p.541—542 / PDF p.33—34}
\end{dataflow}

\begin{crosscheck}
上一条的关键推论是：\textbf{谱上看不见的铝不等于不存在的铝}。因此任何"骨架铝含量"或
"B 酸位数目"的结论都必须回查三件事：（i）同一批样品的 $^{27}$Al MAS NMR 是否同时记录了
四配位、五配位、六配位铝，并说明定量用的是直接极化还是高功率去偶；
（ii）$^{1}$H MAS NMR 的桥羟基积分是否在全弛豫条件下获得（见 8.3.3）；
（iii）脱水／焙烧温度是否与比较样品一致，因为脱羟基程度本身就是温度的函数。三者缺一，
"酸量随温度下降"这类结论就可能只是"铝变不可观测"的假象。
\end{crosscheck}

\subsubsection{酸强度测定：探针分子的低场位移标尺}
\label{sec:acid-strength}

\textbf{强碱探针：吡啶。}强碱性分子氘代吡啶吸附后能有效区分酸性 OH 与非酸性 OH：吡啶与
Si—OH 形成氢键络合物时，其 $^{1}$H 化学位移由 2.0 ppm 移到约 10 ppm；而与桥羟基作用
生成质子化吡啶（吡啶阳离子）时移向更低场，落在 $15.5\sim19.5$ ppm 之间。
\src{S2 §8.3.2，书内 p.542 / PDF p.34}

\textbf{弱碱探针：位移本身就是酸强度的刻度。}不同碱性的吸附分子使桥羟基产生不同的低场
位移。在 \qty{123}{K} 低温下，Haw 等观察到乙烯吸附在 HZSM-5 上使桥羟基向低场位移
2.7 ppm，相同条件下 CO 与乙烷分别只位移 1.8 ppm 与 0.6 ppm。
\src{S2 §8.3.2，书内 p.542 / PDF p.34}

\paragraph{丙酮-2-$^{13}$C：把 B 酸强度做成一条连续刻度。}
Biaglow 等在 \qty{125}{K} 下测定了丙酮-2-$^{13}$C 吸附在多种分子筛上的 $^{13}$C MAS NMR 谱。
在\textbf{每个桥羟基吸附一个丙酮分子}的条件下，相对于纯丙酮分子，羰基 $^{13}$C 的低场位移为：
SAPO-5 上 10.1 ppm，HY 上 12 ppm，HZSM-5 上 15 ppm，HZSM-22 上 18.7 ppm；在 \ce{AlCl3} 与
100\% \ce{H2SO4} 上位移 37 ppm，在 \ce{SbF5} 上位移 42 ppm。\textbf{酸性越强，羰基的 $^{13}$C
化学位移越移向低场}，因此这一方法可把分子筛 B 酸位的强度放在一条从弱到超强的连续刻度上
（图~\ref{fig:acetone-shift}）。\src{S2 §8.3.2，书内 p.542 / PDF p.34}

\begin{figure}[htbp]
\centering
\begin{tikzpicture}
\begin{axis}[
  width=\onecolwidth, height=52mm,
  xbar, bar width=2.4mm,
  xmin=0, xmax=48,
  xlabel={羰基 $^{13}$C 低场位移 $\Delta\delta$ / ppm},
  xlabel style={font=\fontsize{8}{9.5}\selectfont},
  symbolic y coords={SAPO-5,HY,HZSM-5,HZSM-22,AlCl3/H2SO4,SbF5},
  ytick=data,
  y tick label style={font=\fontsize{7.5}{9}\selectfont},
  x tick label style={font=\fontsize{7.5}{9}\selectfont},
  nodes near coords,
  nodes near coords style={font=\fontsize{7}{8.5}\selectfont},
  enlarge y limits=0.10,
  axis lines=left,
]
\addplot[fill=black!30, draw=black] coordinates {
  (10.1,SAPO-5) (12,HY) (15,HZSM-5) (18.7,HZSM-22)
  (37,AlCl3/H2SO4) (42,SbF5)};
\end{axis}
\end{tikzpicture}
\caption{吸附丙酮-2-$^{13}$C 的羰基 $^{13}$C 低场位移，即一条 B 酸强度刻度。读者应看两点：
纵轴自下而上酸强度递增，位移量单调增大；分子筛区间（10.1—18.7 ppm）只占全刻度的前一半，
\ce{AlCl3}／100\% \ce{H2SO4} 与 \ce{SbF5} 把刻度延伸到 42 ppm。条件与局限：\qty{125}{K}、
每桥羟基覆盖一个丙酮分子，位移以纯丙酮为参考态；柱子为原书数据的重绘，未含误差棒。
\redrawmark\figsrc{据 S2 §8.3.2，书内 p.542 / PDF p.34 的数据改绘。}}
\label{fig:acetone-shift}
\end{figure}

\textbf{膦类探针。}吸附三甲基膦（trimethyl phosphine，TMP）与三甲基氧膦（trimethyl
phosphine oxide，TMPO）的 $^{31}$P MAS NMR 化学位移对 B 酸强度同样敏感：
\textbf{B 酸强度越强，$^{31}$P 化学位移越移向低场}。$^{31}$P 的天然丰度 100\%、磁旋比高，
不需要同位素富集，也不必用 CP/MAS，这一点对定量尤其重要（见 8.3.3）。
\src{S2 §8.3.2，书内 p.542 / PDF p.34}

\paragraph{用探针尺寸换取空间信息：全氟丁胺只测外表面。}
把探针分子做大，可以让它\textbf{只能}接触外表面。全氟丁胺
［\ce{(n-C4F9)3N}］的分子直径 0.94 nm，远大于 ZSM-5 的孔径 0.55 nm，因此只能吸附在分子筛
外表面。吸附后 $^{1}$H MAS NMR 谱中 3.9 ppm 处的 B 酸峰强度明显降低，而 6.0 ppm 处的
峰强有所增加——后者来自外表面 B 酸被质子化后向低场位移。由此可以\textbf{通过定量拟合吸附
前后 3.9 ppm 峰的积分面积，算出外表面的 B 酸含量}，同法可算外表面非骨架铝的含量。该实验
还发现吸附全氟丁胺后硅羟基一律向低场位移 $0.2\sim0.3$ ppm，据此推断\textbf{大部分硅羟基位于
分子筛外表面}。\src{S2 §8.3.2，书内 p.542 / PDF p.34}

\begin{dataflow}
上述实验的逻辑链值得单独记住，因为它把"化学位移"与"空间位置"接了起来：

尺寸排斥（探针 0.94 nm $>$ 孔口 0.55 nm）$\rightarrow$ 探针只能到达外表面 $\rightarrow$
只有外表面酸位被质子化 $\rightarrow$ 3.9 ppm 峰\textbf{面积减少量} $=$ 外表面 B 酸量 $\rightarrow$
与总量相减得内表面 B 酸量。注意最后一个箭头依赖一个前提：\textbf{3.9 ppm 峰的减少全部来自
外表面}。若探针能部分进入孔口（例如分子筛有缺陷、介孔或堆积孔），这个前提就被破坏，得到
的"外表面酸量"会被高估。\src{S2 §8.3.2，书内 p.542—543 / PDF p.34—35}
\end{dataflow}

\begin{pitfall}
\textbf{把化学位移直接当作酸强度。}这是本小节最常见、也最难自查的误判，至少有四个来源：
（i）\textbf{参考态漂移}——所有位移都是相对于某个参考态报出的（纯丙酮、液态吡啶、气相
Xe、液态 TMP），换参考态数值整体平移，而排序不变；（ii）\textbf{覆盖度与共吸附物}——丙酮
的位移刻度必须在"每个桥羟基一个丙酮分子"下读取，覆盖度变化或体系残留水都会改变氢键网络，
使位移偏离标定曲线；（iii）\textbf{温度}——弱碱探针的位移刻度依赖低温（\qty{123}{K}、
\qty{125}{K}）冻结交换，室温测量会把不同酸位的信号平均成一条线；（iv）\textbf{孔道限域效应}
——同样的酸位在窄孔中与被质子化探针的构象、取向受约束，$^{31}$P 位移还包含探针本身形变
的贡献。正确做法是：\textbf{在同一次实验、同一参考态、同一温度下测一系列酸强度已知的样品，
用相对次序说话}，并把探针覆盖度写进结论的条件里。
\end{pitfall}

\subsubsection{区分与计数 B 酸／L 酸中心：探针选择与定量纪律}
\label{sec:bl-quant}

\textbf{L 酸只能靠探针。}用 NMR 研究 L 酸性质时没有"直接法"，必须引入探针分子，而且大部分
探针需要同位素标记；好在一般探针既能研究 L 酸，也能同时区分 B 酸。
\src{S2 §8.3.3，书内 p.543 / PDF p.35}

\textbf{\ce{^{13}CO}：L 酸中心的探针，但读不直接。}吸附 CO 的 $^{13}$C 化学位移与其气态相差
很大，可用来探测固体表面的 L 酸中心；然而无论在室温还是低温下，吸附在 L 酸上的 CO 与
物理吸附的 CO、吸附在 B 酸位上的 CO 之间\textbf{都存在交换}，因此谱图上无法直接读出 L 酸位
上 CO 的化学位移，需要\textbf{通过计算}才能获得；该位移落在 $300\sim400$ ppm 之间。
\src{S2 §8.3.3，书内 p.543 / PDF p.35}

\textbf{\ce{^{15}N}-吡啶：位移范围大，但要富集。}Maciel 等对无定形硅铝上 \ce{^{15}N}-吡啶的
\ce{^{15}N} CP/MAS NMR 谱作如下归属：B 酸位上吸附的吡啶在约 205 ppm；L 酸位上的约在
260 ppm；吸附在弱酸位羟基上形成氢键的在 $290\sim300$ ppm；液态吡啶在 315 ppm。$^{15}$N 的
优点是化学位移范围大、各类酸中心容易分开，缺点是天然丰度低，必须同位素富集。
\src{S2 §8.3.3，书内 p.543 / PDF p.35}
\emph{核对说明}：已对照原页面图像（书内 p.543 / PDF p.35），该组数值逐字一致——B 酸位上
吸附吡啶的 \ce{^{15}N} 在 205 ppm 左右、L 酸位上的约在 260 ppm 处、弱酸位羟基上形成氢键的
在 $290\sim300$ ppm、液态吡啶在 315 ppm；原书\textbf{未随数值给出参考标准}（液氨、硝基甲烷等），
跨文献比较前须自行确认参考态。

\paragraph{三甲基膦 TMP：一个探针同时给出 B 酸与 L 酸两个窗口。}
Lunsford 等首先研究了 TMP 在 HY 分子筛上的 $^{31}$P MAS NMR 谱，得到三个分区
（\src{S2 §8.3.3，书内 p.543—544 / PDF p.35—36}）：

\begin{itemize}[leftmargin=1.6em,itemsep=1pt,topsep=2pt]
\item $-3\sim-1$ ppm：TMP 与 B 酸位形成的 \ce{(CH3)3PH+}；
\item $-60\sim-32$ ppm：TMP 与 L 酸位形成的络合物；
\item 约 $-68$ ppm：物理吸附的 TMP。
\end{itemize}

\emph{核对说明}：已对照原页面图像（书内 p.543 / PDF p.35），原书三处数值\textbf{均带负号}
印作 $-3\sim-1$ ppm（\ce{(CH3)3PH+}）、$-60\sim-32$ ppm（TMP–L 酸络合物）、$-68$ ppm 左右
（物理吸附 TMP），与本分片一致；源 OCR 文本层丢失了这些负号，此前据物理次序的恢复方向正确。

这一分区方式带来两个后果，一正一反：\textbf{B 酸位的膦峰远离其他膦峰，因此可以直接测定
表面 B 酸中心的数目}；而\textbf{L 酸络合物峰与物理吸附峰彼此很靠近，故不易直接定量 L 酸的
中心数目}。Coster 等观察到 TMP 在超强酸 \ce{ZrO2}/\ce{SO4^{2-}} 上吸附时，除与 B 酸和 L 酸位
结合的共振峰外，$26$ ppm 处还出现一个强信号；与吸附苯胺后的 ESR 结果比较后，认为该峰
应为 TMP 与\textbf{超强 L 酸中心}的络合物。\emph{核对说明}：已对照原页面图像（书内 p.544 /
PDF p.36），原书此处印作「在 26 ppm 处还出现一个强信号」，\textbf{不带负号}，与本分片
$26$ ppm 的写法一致；与此前 $-60\sim-32$ ppm 的 L 酸络合物窗口相比，该峰更靠低场，
正与「超强 L 酸中心去屏蔽更强」相符。

\begin{figure}[htbp]
\centering
\begin{tikzpicture}
\tikzset{declare function={
  lor(\x,\p,\w) = 1/(1+((\x-\p)/\w)^2);
  spA(\x) = 1.00*lor(\x,-4.6,0.75) + 0.90*lor(\x,-58,2.0) + 0.40*lor(\x,-68,2.5);
  spB(\x) = 1.00*lor(\x,-2.4,0.85) + 0.90*lor(\x,-45,2.1) + 0.40*lor(\x,-68,2.5);
  spC(\x) = 1.00*lor(\x,-1.0,0.95) + 0.90*lor(\x,-34,2.2) + 0.40*lor(\x,-68,2.5);
}}
\begin{axis}[
  width=\onecolwidth, height=56mm,
  xmin=-76, xmax=6, x dir=reverse,
  ymin=0, ymax=4.15,
  xlabel={$\delta(^{31}\mathrm{P})$ / ppm（低场 $\leftarrow$ $\rightarrow$ 高场）},
  xlabel style={font=\fontsize{8}{9.5}\selectfont},
  ylabel={相对强度（叠加显示）},
  ylabel style={font=\fontsize{8}{9.5}\selectfont},
  axis lines=left, ytick=\empty,
  xtick={2,0,-10,...,-70},
  tick label style={font=\fontsize{7.5}{9}\selectfont},
  legend style={font=\fontsize{7.5}{9}\selectfont, draw=none, fill=none,
                legend columns=3, at={(0.5,1.02)}, anchor=south},
  samples=600, domain=-76:6,
]
\addplot[black, thick] {spA(x)};
\addplot[black, thick, densely dashed] {spB(x)+1.35};
\addplot[black, thick, dotted] {spC(x)+2.70};
\legend{弱酸位, 中等酸位, 强酸位}
\node[anchor=north west, align=left, text width=20mm,
      font=\fontsize{7.5}{9}\selectfont] at (axis cs:6,4.12)
      {B 酸位 \ce{(CH3)3PH+}\\ $-3\sim-1$ ppm};
\node[anchor=north east, align=right, text width=20mm,
      font=\fontsize{7.5}{9}\selectfont] at (axis cs:-76,4.12)
      {L 酸位络合物\\ $-60\sim-32$ ppm};
\end{axis}
\end{tikzpicture}
\caption{模型计算的吸附 TMP 的 $^{31}$P MAS NMR 谱：两条峰位随酸强度移动。读者应看两件事：
（i）每个酸强度下都有两组信号——左（低场）为 B 酸位上的 \ce{(CH3)3PH+}，右（高场）为 L 酸
位络合物；（ii）酸强度由弱到强（自下而上），\textbf{两组峰同时向低场移动}，B 酸峰由
$-4.6$ 移到 $-1.0$ ppm，L 酸峰由 $-58$ 移到 $-34$ ppm，而 $\approx-68$ ppm 的物理吸附峰
（图中最右的小峰）位置不随酸强度改变。条件与局限：曲线由洛伦兹线型叠加合成，峰位取自
原书给出的分区端点，线宽、相对强度与叠加偏置均为设定值，\textbf{不是实验谱}；真实谱中 L 酸
络合物峰与物理吸附峰相距很近、部分重叠，定量须做分峰拟合。\modelmark\figsrc{模型参数据
S2 §8.3.3，书内 p.543—544 / PDF p.35—36 的化学位移区间设定。}}
\label{fig:tmp-p}
\end{figure}

\begin{table}[htbp]
\centering
\footnotesize
\setlength{\tabcolsep}{1.5pt}
% 原五列合计仅 80 mm（版心 165 mm），最窄列装不下内容，单元格溢出 9.05 pt。
% 【重叠修正，2026-09-30】上一轮只加宽到 108 mm，仍空着 57 mm 版心；单元格被挤成
% 多行，表格连图题一起高出整页 —— 编译日志实测
%   `LaTeX Warning: Float too large for page by 37.14322pt on input line 361`，
% 浮动体被强行放下后**溢出到页脚**：图题最后两行落在 y=788—809，
% 而页脚页码 `158` 在 y=784.7—791.0，实测 5 对字形墨迹相交（最大 `表`×`5` 0.907 pt²）。
% 本轮把五列吃满版心（158 mm + 4×2×1.5 pt 列间距 ≈ 162 mm < 165 mm），
% 单元格行数随之下降，浮动体不再超高。列宽之和勿减，否则会再次触发上述警告。
\begin{tabular}{@{}p{20mm}p{12mm}p{30mm}p{40mm}p{56mm}@{}}
\toprule
探针分子 & 核种 & 测什么 & $\delta$ 区间 / ppm & 定量前提与主要干扰 \\
\midrule
\ce{(CH3)3P}（TMP）
 & \ce{^{31}P} & 同时给出 B 酸位与 L 酸位两个窗口 &
 B：$-3\sim-1$；\\ L：$-60\sim-32$；物理吸附 $\approx-68$ &
 丰度 100\%，不需富集、不需 CP/MAS；L 络合物峰与物理吸附峰靠近，须分峰 \\
\addlinespace[2pt]
\ce{^{15}N}-吡啶
 & \ce{^{15}N} & 区分 B 酸／L 酸／弱酸位 &
 B 205；L 260；\\ 弱酸 $290\sim300$；液态 315 &
 位移范围大、易分辨；天然丰度低，\textbf{必须同位素富集}，参考态须统一 \\
\addlinespace[2pt]
丙酮-2-$^{13}$C
 & \ce{^{13}C} & B 酸强度刻度 &
 相对纯丙酮的位移 $10.1\sim42$ &
 须"每桥羟基一个丙酮分子"、低温（\qty{125}{K}）；覆盖度与残水改变刻度 \\
\addlinespace[2pt]
\ce{^{13}CO}
 & \ce{^{13}C} & L 酸中心 &
 L 酸位 $300\sim400$（须计算求得） &
 室温与低温下 L 酸位、B 酸位、物理吸附的 CO 均在交换，谱上读不出，只能拟合 \\
\addlinespace[2pt]
吡啶-d$_5$
 & $^{1}$H & 区分酸性／非酸性 OH &
 氢键 $\approx10$；\\ 吡啶阳离子 $15.5\sim19.5$ &
 强碱优先占据最强酸位；须确认化学计量覆盖，否则高估强酸位数 \\
\addlinespace[2pt]
乙烯／CO／乙烷
 & $^{1}$H & 弱碱探针测桥羟基低场位移 &
 $+2.7$／$+1.8$／$+0.6$（\qty{123}{K}） &
 位移小，必须低温冻结交换；参考态误差与位移同量级，只宜作相对比较 \\
\addlinespace[2pt]
全氟丁胺
 & $^{1}$H & 外表面 B 酸与非骨架铝定量 &
 3.9 峰降、\\ 6.0 峰升（位移 $0.2\sim0.3$） &
 分子 0.94 nm $>$ 孔口 0.55 nm，只接触外表面；靠差减积分，前提是无介孔／缺陷漏入 \\
\addlinespace[2pt]
$^{1}$H（直接法）
 & $^{1}$H & 各类羟基的归属与桥羟基定量 &
 分六档，$0.5\sim7.0$ &
 须真空脱水、高速 MAS（或 CRAMPS）；残水与装样污染同时改变位移与强度 \\
\bottomrule
\end{tabular}
\caption{酸性探针分子对照。读者用法：先按"想测什么"选行，再核对同一行的"定量前提与
主要干扰"是否满足；表中位移均为相对各自参考态的数值，\textbf{不可跨探针直接比大小}。
条件与局限：区间取原书给出的端点，个位数与符号见正文各处的待核对标注。
\src{据 S2 §8.3.2—§8.3.3，书内 p.542—544 / PDF p.34—36 汇总。}}
\label{tab:probes}
\end{table}

\begin{principle}
\textbf{定量前提：小扳倒角 $+$ 足够长的弛豫延迟。}探针定量的目标是让\textbf{峰面积比等于
核数目比}。单脉冲实验中第 $n$ 次扫描的稳态信号为

\begin{equation}
\label{eq:ernst}
S\;\propto\;N\cdot\frac{\sin\alpha\,\bigl(1-e^{-D_1/T_1}\bigr)}
{1-\cos\alpha\,e^{-D_1/T_1}}
\end{equation}

式中 $N$ 为产生该峰的核数目，$\alpha$ 为扳倒角（射频脉冲宽度决定），$D_1$ 为弛豫延迟，
$T_1$ 为自旋--晶格弛豫时间。式~\eqref{eq:ernst} 说明两件事：（i）只要
$D_1\gtrsim5T_1$，指数项趋于零，$S\propto N\sin\alpha$，面积比即数目比——\textbf{弛豫延迟
不足时，$T_1$ 短的峰被相对放大，强度比失真}；（ii）当 $\alpha$ 小到
$\sin\alpha\approx\alpha$、$\cos\alpha\approx1$ 时，$D_1/T_1$ 项整体约掉，
$S\propto N\alpha$，\textbf{强度比对 $T_1$ 的差异不再敏感}。这就是"小扳倒角
（$\alpha\lesssim10°$）$+$ 强射频脉冲"的用处：它把定量从"必须等 5$T_1$"变成"不必等 5$T_1$"，
代价是灵敏度按 $\sin\alpha$ 下降，需要更多累加次数。\pedset{本式为 NMR 定量的通用关系，
是教学补充（原书未给该式）；原书给出的要求是定性表述——固体 NMR "在定量方面优于红外光谱"，
但 L 酸络合物峰与物理吸附峰靠近而"不易直接定量"（\src{S2 §8.3.3，书内 p.543—544 /
PDF p.35—36}）。}
\end{principle}

\begin{measure}
\textbf{做定量前必须逐条确认的六件事。}（i）\textbf{扳倒角}：改用 $\alpha\lesssim10°$ 的
小扳倒角（或用已知 $T_1$ 的模型化合物标定脉冲宽度），并保留强射频脉冲以覆盖全谱宽；
（ii）\textbf{弛豫延迟}：实测各峰的 $T_1$ 后取 $D_1\gtrsim5T_1$，$\alpha$ 很小且已验证稳态时
可按 $(1-\cos\alpha)T_1$ 量级放宽；（iii）\textbf{旋转边带}：把转速提高或把全部边带一并
积分，只积中心带会系统性低估；（iv）\textbf{探针化学计量}：确认吸附量对应"每酸位一个
探针分子"，过量探针会同时产生物理吸附峰；（v）\textbf{温度与交换}：弱碱探针必须在低温
（\qty{123}{K}、\qty{125}{K}）测，否则不同酸位的信号被交换平均；（vi）\textbf{分峰模型}：
L 酸络合物与物理吸附峰重叠时应报告拟合约束与误差，而不是给一个单点数值。
\end{measure}

\begin{pitfall}
\textbf{四个必须设防的误判。}
（1）\textbf{弛豫延迟不足导致定量错误。}这是本节最重要的判断纪律：探针分子与酸位结合后
运动性下降，$T_1$ 与物理吸附态、自由态可以相差一个数量级；用统一的短延迟采集，会把
$T_1$ 短的物种（常是物理吸附或弱结合态）算多，把强结合态算少，而"B 酸／L 酸数目比"恰恰
建立在这两类峰的面积比上。看到一篇文献给酸位数目而不给 $\alpha$、$D_1$、$T_1$，
就应当把它当作\textbf{定性}结果使用。
（2）\textbf{把化学位移直接当酸强度。}位移受参考态、覆盖度、温度与限域效应共同影响
（见 8.3.2 的误判框），只能在同一套标定条件下作相对比较。
（3）\textbf{探针分子太大进不去小孔，被误读为"没有酸位"。}探针尺寸既是优点也是陷阱：
0.94 nm 的全氟丁胺进不去 0.55 nm 的 ZSM-5 孔道，那是有意为之；但若用大尺寸探针去筛
小孔分子筛并据此宣布"孔内无酸性位"，就把\textbf{扩散限制}错当成了\textbf{无酸位}。
判据是：换一个尺寸小于孔口、化学性质相近的探针（如 TMP、\ce{NH3}）重测，若信号出现，
则原结论只是可接近性问题。分子筛中扩散慢是常态——源文在 \ce{^{129}Xe} 一节记录到，苯在
NaY／NaX 中室温扩散很慢，吸附 120 天后才出现多个分离的共振峰；\qty{523}{K} 加热 10 h 才
给出均匀分布的单峰（\src{S2 §8.4.2，书内 p.549—550 / PDF p.41—42}）。装样后的
\textbf{平衡时间}因此是探针实验的隐藏变量，必须与升温／抽空历史一并报告。
（4）\textbf{把旋转边带当成真实峰。}MAS 下只要转速不足以平均该核的各向异性，中心带两侧就
会出现间隔\textbf{恰好等于转速}的旋转边带；$^{31}$P 与 \ce{^{129}Xe} 的化学位移跨度大，
边带更容易落到"看起来像另一个酸位"的位置上，从而凭空多出一个"L 酸位"。判别办法只有一条：
\textbf{改变转速重测}——真实峰的位置不随转速移动，边带的位置随转速整体平移；此外可用
边带间隔直接反读转速自检。\src{S2 §8.3—§8.4，书内 p.543—546 / PDF p.35—38}
\end{pitfall}

\begin{crosscheck}
\textbf{酸性结论的交叉验证组合。}（i）$^{1}$H MAS NMR 的桥羟基积分（B 酸总量）与
$^{27}$Al MAS NMR 的四配位骨架铝含量互相校核，两者不等时先怀疑"铝变不可观测"，而不是
直接修改酸量结论；（ii）吡啶吸附红外（区分 B／L 酸的特征环振动）与 $^{15}$N-吡啶 NMR 互相
校核——红外给出类型与相对强度、NMR 给出可分辨的位移与数目；（iii）\ce{NH3}-TPD 的脱附峰温
给出强度分布的定性排序，与丙酮-$^{13}$C 的位移刻度对照，若排序矛盾，优先检查探针覆盖度与
平衡时间；（iv）外表面酸量的结论必须用不同尺寸探针（全氟丁胺 vs TMP）分别测量并给出
差值；（v）所有定量结论最终要写清 $\alpha$、$D_1$、转速、探针覆盖度与平衡时间四项
实验参数，否则不可复现。\src{S2 §8.3，书内 p.540—544 / PDF p.32—36}
\end{crosscheck}

源文对 NMR 测酸性的总体评价是：固体 NMR 已成为研究多相催化剂酸性的有力工具，能得到其他
方法难以得到的信息，\textbf{在定量方面优于红外光谱}；但在 L 酸研究方面一般需采用同位素
富集的探针分子，相比之下\textbf{三甲基膦不失为一个较好的 NMR 探针}。
\src{S2 §8.3.3，书内 p.544 / PDF p.36}

\subsection{催化剂表面吸附分子的 NMR 研究}
\label{sec:adsorbed}

研究催化剂表面吸附分子的化学性质以及催化反应中间产物的结构，是从分子水平阐明催化反应
机理的关键，NMR 已被广泛用于表面吸附分子状态与氢溢流的研究；多相催化过程在很大程度上
取决于反应物在催化剂表面的吸附行为，而惰性气体氙（Xe）是一种非常理想的探针分子，
\ce{^{129}Xe} NMR 已成为研究多孔物质、特别是分子筛结构的有效工具。
\src{S2 §8.4，书内 p.544 / PDF p.36}

\subsubsection{分子筛晶体孔道中吸附有机物的化学状态}
\label{sec:methanol}

\textbf{案例：甲醇在 HZSM-5 上的吸附与转化。}甲醇转化制烃是工业上有代表性的重要反应，
其机理（特别是第一个 C—C 键如何生成、甲醇与骨架如何相互作用）长期有争议。用
$^{1}$H MAS NMR 研究甲醇在 HZSM-5 上的吸附性质，得到的观察与推断可以分开陈述。
\src{S2 §8.4.1，书内 p.544—545 / PDF p.36—37}

\begin{obs}
甲醇吸附在 HZSM-5 上后，$^{1}$H MAS NMR 谱中出现 4.1 ppm 的甲基峰，OH 峰位于
9.1 ppm，比甲醇在 \ce{CDCl3} 中向低场位移 3.1 ppm；不论吸附的是 \ce{CD3OH}、\ce{CH3OD}
还是 \ce{CD3OD}，所有羟基的化学位移都在 9.1 ppm 左右。把覆盖度降到每个 B 酸位"只覆盖
两个甲醇分子"时，OH 信号由 9.1 ppm 移向 10.5 ppm。
\emph{核对说明}：已对照原页面图像（书内 p.544 / PDF p.36），4.1 ppm 甲基峰、9.1 ppm OH
峰、「比其在 \ce{CDCl3} 中向低场位移 3.1 ppm」、三种氘代样品羟基均在 9.1 ppm 左右、以及
「每个 B 酸只覆盖两个甲醇分子时 OH 由 9.1 ppm 移向 10.5 ppm」均逐字一致。两点说明：
原书\textbf{未印出} \ce{CDCl3} 中甲醇 OH 的具体位移，约 6.0 ppm 系由 $9.1-3.1$ 反推的本分片
推算值；原书原句为「当每个 B 酸\textbf{只}覆盖两个甲醇分子时」，对应的是\textbf{低覆盖}一侧。
\end{obs}

\begin{infer}
液体甲醇中 OH 之间已存在氢键，O—H 键的静电极化使其 $^{1}$H 去屏蔽；吸附到 HZSM-5 上后
出现如此大的低场位移，\textbf{必定来自很强的氢键和／或醇的质子化}。三种同位素取代样品的
OH 位移都在 9.1 ppm 附近，说明\textbf{所有 $^{1}$H 在 NMR 时域上是等价的}：在甲醇形成的
氢键中，端基甲醇分子与 B 酸位的桥氧形成氢键，每个甲醇分子中的甲氧基离子相同，这样一个
以氢键相连的带电簇可在孔道中转动，使羟基的 H 变成等价。覆盖度降低时 OH 由 9.1 移到
10.5 ppm，说明 B 酸位的 H 与吸附甲醇的 OH 之间发生了\textbf{快速交换}。
\src{S2 §8.4.1，书内 p.544—545 / PDF p.36—37}
\end{infer}

由上述结果得到一条可直接使用的判据：\textbf{用吸附甲醇中 OH 的低场位移可以测量固体酸
催化剂的供质子能力，低场位移越大，B 酸位酸性越强。}HY 与 HL 分子筛上甲醇 OH 的位移小于
HZSM-5，与它们酸性较弱相符；同 Si/Al 比而合成方法不同的 HZSM-5，吸附甲醇的 OH 低场位移
也不同——说明位移对分子筛的\textbf{类型与制备历史}都敏感。钠型分子筛是反例：NaZSM-5 上
吸附甲醇时 OH 峰比液体甲醇向\textbf{高场}位移 1 ppm，因为甲醇通过氧与 Na$^{+}$ 配位、
打断了氢键；对 NaZSM-5、NaY、NaA，吸附甲醇的 OH 低场位移依次增加，说明甲醇与骨架形成
氢键的可能性依次增加。\src{S2 §8.4.1，书内 p.545 / PDF p.37}

\textbf{$^{13}$C MAS NMR 跟踪反应：从弱键分子到强键中间物。}原书图 8-12 给出甲醇吸附在
SAPO-5 上的 $^{13}$C MAS NMR 谱（(1) 室温；(2) \qty{150}{\degC} 加热 10 min；
(3) \qty{250}{\degC} 加热 10 min）。谱学事实与归属如下
（\src{S2 §8.4.1，书内 p.545—546 / PDF p.37—38}）：室温时在 50 ppm 出现一个窄峰，归属为
相对高运动性的甲醇分子；\qty{150}{\degC} 加热 10 min 后有 37\% 的甲醇转化为二甲醚，峰形
仍然较窄，说明分子仍有很大的运动性；继续 \qty{250}{\degC} 加热 10 min 后谱峰加宽，呈典型
的固体谱，除甲醇峰外出现一个新峰，其化学位移与二甲醚相近，归属为与骨架相连的甲氧基
（\ce{CH3-O-Si}）反应中间物；再加热到 \qty{300}{\degC}，甲醇转化为烯烃与烷烃的混合物，
谱线又变窄，因为产物的流动性增强。

\begin{dataflow}
把上段读成一条"键强度"的轨迹：\textbf{弱键的甲醇分子（窄线）$\rightarrow$ 强键的反应中间物
（宽线，甲氧基）$\rightarrow$ 弱键的碳氢产物（窄线）}。它的判据是\textbf{线宽}：窄线对应大
运动性、宽线对应被骨架固定。由此得到一个重要的方法论结论：吸附有机分子的 $^{1}$H 与
$^{13}$C MAS NMR 研究，可为确定各类分子筛的反应中间物、探索反应机理提供十分重要的证据。
\src{S2 §8.4.1，书内 p.546 / PDF p.38}
\end{dataflow}

\begin{pitfall}
\textbf{线宽与化学位移不是"结构",而是"结构与运动的合成"。}（i）在变温实验中，峰变宽既可
来自生成了强键中间物，也可能只是运动性下降或交换速率落入中间窗（$k\approx\Delta\nu$），
两者只能靠\textbf{变温序列 $+$ 化学位移归属}一起区分；（ii）温度与加热时间是结论的一部分
——上例的"37\% 转化"是 \qty{150}{\degC}、10 min 这个条件下的数值，换条件就换数值；
（iii）归属"与二甲醚位移相近"的峰时，必须用独立参照（纯二甲醚吸附谱、甲氧基模型化合物）
回查，不能只凭位移接近就定为甲氧基；（iv）这类实验的峰位随温度漂移，参考态与温度必须
写进结论。
\end{pitfall}

\subsubsection{吸附探针分子 \texorpdfstring{\ce{^{129}Xe}}{129Xe} 的 NMR 研究}
\label{sec:xe}

\ce{^{129}Xe} 是理想探针：核自旋量子数 $I=1/2$（无四极增宽）、天然丰度 26.4\%、共振频率比
$^{13}$C 高约 10\%；若不考虑弛豫时间，其检测灵敏度是 $^{13}$C 的 31.8 倍。单原子
\ce{^{129}Xe} 的 $T_1$ 很长，但吸附到分子筛后受顺磁物种与痕量杂质影响 $T_1$ 大幅降低，
使 \ce{^{129}Xe} NMR 成为一种方便的、灵敏度高于 $^{13}$C 的方法。Xe 核的球形电子云很大，
对周围环境十分敏感：分子筛的结构、组成与孔道的任何微小变化都会改变 Xe 的电子云密度，
从而改变共振峰的化学位移。\src{S2 §8.4.2，书内 p.546 / PDF p.38}

\begin{principle}
\textbf{吸附 Xe 的化学位移是几项贡献之和}，每一项对应一种不同的干扰
（\src{S2 §8.4.2，书内 p.546 / PDF p.38}）：

\begin{equation}
\label{eq:xe-shift}
\delta=\delta_{0}+\delta_{E}+\delta_{M}+\delta_{\mathrm{Xe\text{-}Xe}}+\delta_{S}+\delta_{\mathrm{SAS}}
\end{equation}

式中 $\delta_{0}$ 为 Xe 气体在吸附压力外推为零时的化学位移，取 $\delta_{0}=0$ 作为化学位移的
参考；$\delta_{E}$、$\delta_{M}$ 为分子筛中阳离子的电场与磁场贡献，对小离子（Na$^{+}$、
Li$^{+}$、H$^{+}$ 等）其 $\delta_{M}=0$，室温下 $\delta_{E}$ 可忽略；$\delta_{\mathrm{Xe\text{-}Xe}}$
为 Xe—Xe 相互作用的贡献，与 Xe 的吸附浓度有关；$\delta_{S}$ 为孔道硅铝壁与氙相互作用的
贡献，直接与孔道结构有关；$\delta_{\mathrm{SAS}}$ 为孔内强吸附位的贡献。Xe 探针提供的
信息一般由 $\delta=f(n)$ 曲线得到，$n$ 为每克无水固体吸附的 Xe 原子数。

\textbf{把 $\delta_{S}$ 与孔尺寸连起来。}为从化学位移得到未知结构分子筛的自由空间或结构
缺欠，用平均自由程 $l$ 把 $\delta_{S}$ 与内部体积的大小和形状关联：$l$ 是一个 Xe 原子在自由
空间无规运动时两次连续碰撞之间运动距离的平均值。Fraissard 等提出

\begin{equation}
\label{eq:fraissard}
\frac{1}{\delta_{S}}=\frac{1+l/a}{\delta_{a}}
\end{equation}

其中 $\delta_{a}$ 与分子筛结构有关、因而也与 $l$ 有关；$l$ 直接与内部体积的大小和形状有关，
于是可由 $\delta_{S}$ 的大小分析孔的几何性质。
\emph{核对说明}：已对照原页面图像（书内 p.547 / PDF p.39），原书以行内式印作
$1/\delta_{S}=(1+\bar{l}/a)/\delta_{a}$，结构和本分片式~\eqref{eq:fraissard} 一致；差别仅在原书
用带上划线的 $\bar{l}$ 表示平均自由程（本分片统一写作 $l$）。原书对 $a$ 未作任何说明，对
$\delta_{a}$ 也只说「与分子筛结构有关，因此也与 $\bar{l}$ 有关」，两者在正文中\textbf{均无完整定义}
——本分片保留这一缺口，不代为补定义。
\end{principle}

\begin{measure}
\textbf{怎么测：把化学位移写成吸附量的函数。}标准做法是测定 $\delta=f(n)$ 曲线，$n$ 为每克
无水固体的 Xe 原子数，并把曲线外推到 Xe 浓度为零，取 $\delta(n\to0)$ 作为表征孔结构的量。
判读规则（\src{S2 §8.4.2，书内 p.547 / PDF p.39}）：（i）只有一种孔结构时谱上只出现一条
吸附位单线；（ii）大孔分子筛（如 Y 型、ZK4）中 Xe—Xe 碰撞是各向同性的，$\delta=f(n)$ 为
直线，斜率 $\mathrm{d}\delta/\mathrm{d}n$ 正比于局部 Xe 原子密度、反比于分子筛内部自由体积；
（iii）窄通道分子筛（孔径 $4\sim7.5$ \AA{}）中 Xe—Xe 碰撞是各向异性的，斜率随每个笼或部分
通道吸附的 Xe 原子数而增加；（iv）孔道或笼很小、或 Xe 扩散受更大限制时，$\delta(n\to0)$
变大；（v）小孔分子筛中 Si/Al 值相对更重要——平均自由程不变时 $\delta_{S}$ 也会变化，
必须考虑分子筛表面化学组成的影响。
\end{measure}

\paragraph{孔结构与多重吸附位。}当 Xe 在具有两种类型孔道的分子筛上吸附时，若不同吸附位上
Xe 的交换速率较慢，可检测到两个 \ce{^{129}Xe} 信号。丝光沸石的孔结构由一维通道和与之
垂直的边袋组成：温度低到足以阻止 Xe 在两种吸附位之间交换时，其 \ce{^{129}Xe} 谱出现两个
信号，分别对应通道与边袋，\textbf{两个信号的合并温度取决于阳离子性质}——H 型丝光沸石为
\qty{273}{K}，Na 型为 \qty{370}{K}；若是 Cs$^{+}$，边袋不吸附 Xe，只有一条谱线。
Moudrakovski 等用二维 NMR 方法测定了主通道与边袋之间 Xe 交换的速率常数。当观察到多个
\ce{^{129}Xe} 信号时，在给定总 Xe 浓度下可利用谱线强度随温度和阳离子位置的变化来研究
Xe 在不同自由空间中的分布，由此得到分子筛结构（内部生长与结晶）的有用信息。
\src{S2 §8.4.2，书内 p.547 / PDF p.39}

\begin{dataflow}
\textbf{从 $\delta=f(n)$ 曲线读出四类结论。}
（1）\textbf{结晶度}：在一定压力下原位观测合成过程中的 \ce{^{129}Xe} 谱线，与结晶完好的
标准分子筛比较谱线强度即可确定结晶度；若室温下非结晶相吸附 Xe 的量可忽略，则比较被测
样品与标准分子筛 $\delta=f(n)$ 曲线的斜率，由斜率比也可测量结晶度。某些分子筛处理后会
出现第二种孔性质，谱上表现为新增一条谱线——在 Y 型与丝光沸石中都观察到过。
（2）\textbf{组成}：\ce{^{129}Xe} 化学位移与化学组成（特别是电荷数即 Si/Al 值）有关，低温与
小孔径时依赖更明显。八面沸石类分子筛随 Si/Al 增大（Al 减少）$\delta_{S}$ 成倍减少，但变化
非线性，取决于孔大小与脱铝形成的非骨架铝物种；窄孔道的 ZSM-5、ZSM-11 中
$\delta_{S}$ 与 Al 浓度的依赖关系大于大笼八面沸石，且 $\delta_{S}$ 与 Al 含量\textbf{呈线性
关系}；值得注意的是\textbf{每个晶胞含两个 Al 原子时该直线出现断裂}（原书图 8-13），说明
在这个浓度下 Xe 与孔壁的相互作用 $\delta_{S}$ 发生了变化。据此，\ce{^{129}Xe} NMR 可以完全
区分 ZSM-5 与 ZSM-11 这样结构类似的分子筛。
（3）\textbf{结构类型}：分子筛结构不同，$\delta=f(n)$ 曲线完全不同——ZSM-48 的曲线是斜率
很大的直线，截距约 108 ppm；Nu-3 的曲线则很平缓，低覆盖率时 $\delta$ 几乎不变。后者与其
结构有关：各笼由不对称八元环分开，Xe 通过速率很慢，低覆盖度时 Xe 相隔很远，
$\delta_{\mathrm{Xe\text{-}Xe}}$ 为零，因而 $\delta$ 与吸附量无关。
（4）\textbf{吸附相分布}：见下。
\src{S2 §8.4.2，书内 p.547—548 / PDF p.39—40}
\end{dataflow}

\paragraph{阳离子的影响。}当 Xe 在某些吸附位是强吸附时，必须考虑 $\delta_{\mathrm{SAS}}$ 项：
低吸附压时 $\delta_{\mathrm{SAS}}$ 占优势；压力增加后吸附发生在弱吸附位，化学位移降低，它是
$\delta_{\mathrm{Xe\text{-}Xe}}$ 与 $\delta_{\mathrm{SAS}}$ 的权重平均；压力更高时
$\delta_{\mathrm{Xe\text{-}Xe}}$ 又占主导，$\delta$ 随吸附压增加，因此 $\delta=f(n)$ 曲线出现
一个\textbf{最低点}。在 X 型与 Y 型分子筛中含 Mg$^{2+}$、Ca$^{2+}$、Zn$^{2+}$、稀土离子
Y$^{3+}$、La$^{3+}$、Ce$^{3+}$ 乃至顺磁离子 Ni$^{2+}$、Co$^{2+}$、Ru$^{3+}$ 时都观察到这种
特征；当 Xe 只与超笼中的阳离子相互作用时，可以确定阳离子的位置。室温下 X 型与 Y 型分子筛
中 H$^{+}$ 与 Na$^{+}$ 对 \ce{^{129}Xe} NMR 的影响可忽略，且与 Si/Al 值、H$^{+}$ 或
Na$^{+}$ 的数量无关；$\delta$ 值随 Xe—Xe 相互作用而增加，曲线取决于阳离子类型——\textbf{离子
越重，$\delta$ 越大}，如 KY 与 RbY 的 $\delta(n\to0)$ 分别为 78 ppm 与 99 ppm。
\src{S2 §8.4.2，书内 p.548—549 / PDF p.40—41}

Bonardet 与 Fraissard 研究 MgY 上吸附 Xe 的谱：离子交换度 $\lambda<53\%$ 时 Mg$^{2+}$ 位于
方钠石笼中或六方棱柱中，$\delta$ 的变化与 NaY 一样、与 Xe 浓度呈线性关系；$\lambda>53\%$
时 Mg$^{2+}$ 位于超笼中，对应每一个 $\lambda$ 值 $\delta=f(n)$ 曲线都出现最低点，$\lambda$
越高最低点越移向高浓度，$\delta(n\to0)$ 大约正比于 Mg$^{2+}$ 吸附的 Xe 原子核处\textbf{电场
的平方}。二价阳离子的曲线显示较大的低场位移与抛物线特征，主要由于二价阳离子产生的强
电场引起 Xe 高度极化与电子云畸变。对顺磁离子 Ni$^{2+}$、Co$^{2+}$ 必须考虑 $\delta_{M}$ 项，
它可以使 $\delta$ 达到几千 ppm，此时可研究阳离子的位置与氧化态；对每种离子交换度，都可以
通过改变预处理温度来研究二价阳离子环境的变化及从超笼的迁移。三价阳离子交换的八面沸石
（Y$^{3+}$-NaY、La$^{3+}$-NaY、Ce$^{3+}$-NaX、Y$^{3+}$-NaX）的 $\delta=f(n)$ 曲线也呈带最低点
的抛物线形，表明 Xe 在超笼内三价离子上的吸附仍为强吸附；相反 CeX、LaX、LaY 上的位移表明
Ce$^{3+}$ 与 La$^{3+}$ 位于六方棱柱或方钠石笼中，并不吸附 Xe 原子。顺磁离子 Ru$^{3+}$ 引起
很大的低场位移，由其化学位移与线宽变化可知：室温还原时 Ru 颗粒高度分散并位于分子筛笼内，
而真空下发生金属迁移。\src{S2 §8.4.2，书内 p.549 / PDF p.41}

\paragraph{确定吸附相的分布。}Gedeon 等用 \ce{^{129}Xe} NMR 研究 HY 在不同脱水度下水的分布
（以相对水浓度 $C$ 表示，饱和含水 $C=1$）：$\delta(n\to0)$ 随含水量降低而迅速减少，到
$C=0.15$（相当于每个方钠石笼中有四个水分子）后保持不变，说明 $C<0.15$ 时 Xe 的平均自由程
不再变化，超笼与六方棱柱已完全脱水，剩下的水分子位于方钠石笼中；$\delta=f(n)$ 曲线的斜率
在 $C$ 由 1 降到 0.4 的过程中也迅速降低、之后保持不变，说明这一区间内可吸附 Xe 的微孔
自由体积逐渐增加，$C<0.4$ 时保持不变、超笼中的水已完全脱除。因此 $0.15<C<0.4$ 时水分子对
孔有部分堵塞，取决于水化程度；窗口的脱水几乎不影响内部孔径，但使 Xe 更容易在 $\alpha$-笼
之间扩散。\src{S2 §8.4.2，书内 p.549 / PDF p.41}

Liu 等用 \ce{^{129}Xe} NMR 考察 NaY 与 NaX 中吸附苯的分布：室温下苯在孔道内扩散很慢，
吸附 120 天后谱上出现几个共振峰，说明苯分子非均匀地分布在分子筛笼中；样品加热到
\qty{523}{K} 保持 10 h 则出现一个单峰，对应苯的均匀分布。谱线宽度随苯的平均覆盖度 $\theta$
变化很大：$\theta<2.5$ 时线宽随 $\theta$ 线性增加，原因是可接受 Xe 的自由体积减小，或吸附
苯的位间跳跃增加而使平均自由程减小；线宽与 Xe 浓度无关，说明 Xe 可以在超笼之间快速运动，
减少了 Xe—Xe 相互作用。$2.5<\theta<3.5$ 时线宽在 $\theta=3$ 处出现极大值，其位置取决于 Xe
的吸附量，原因是苯—Na 络合物向超笼中心协同迁移引起分子重排，Xe 与苯之间存在相互障碍；
$\theta>4$ 时线宽重新增加，表明苯在超笼有限自由空间内有协同相互作用；$\theta>4.5$ 时苯分子
排列在超笼内，Xe 原子不可能再进入。\src{S2 §8.4.2，书内 p.549—550 / PDF p.41—42}

\begin{pitfall}
\textbf{\ce{^{129}Xe} 谱的五类误判。}（i）\textbf{把 $\delta$ 直接归给一项贡献}：式~\eqref{eq:xe-shift}
中 $\delta_{\mathrm{Xe\text{-}Xe}}$、$\delta_{S}$、$\delta_{\mathrm{SAS}}$ 同时存在，且随压力
此消彼长，只有 $\delta=f(n)$ 曲线（而不单点位移）才可解耦；只看一个压力点的位移就断言
"孔径变小"是典型错误。（ii）\textbf{把两个信号读成两种孔}：两个信号也可能来自交换慢的
同一类孔中的两种吸附位（如丝光沸石的通道与边袋），判别要靠\textbf{变温合并实验}——合并
温度本身就是信息（H 型 273 K、Na 型 370 K）。（iii）\textbf{忽略顺磁杂质}：Ni$^{2+}$、Co$^{2+}$、
Ru$^{3+}$ 使 $\delta_{M}$ 达几千 ppm 并极大加宽谱线，此时结构与阳离子位置的信息被掩盖，
须先确认样品中是否含顺磁物种。（iv）\textbf{把线宽当成单一因素的函数}：线宽同时受自由体积、
位间跳跃、协同重排与 Xe 浓度影响，上例中同一条线宽曲线在 $\theta<2.5$、$\theta=3$、
$\theta>4$ 三段由不同机制主导。（v）\textbf{平衡时间不足}：苯在 NaY 中需要 120 天（室温）
或 \qty{523}{K}/10 h 才达到均匀分布，把未平衡的谱当作平衡态分布，会把\textbf{动力学}差异
读成\textbf{结构}差异。\src{S2 §8.4.2，书内 p.547—550 / PDF p.39—42}
\end{pitfall}

\begin{crosscheck}
\ce{^{129}Xe} NMR 给出的都是\textbf{孔内自由空间与吸附相}的信息，必须与下列手段回查：
（i）$\delta_{S}$／$\delta(n\to0)$ 与结构参数对照——用已知结构的分子筛（Y、ZK4、ZSM-5、
ZSM-11、ZSM-48、Nu-3、丝光沸石）建立 $\delta(n\to0)$、$\mathrm{d}\delta/\mathrm{d}n$ 与孔径／
笼尺寸的经验关系；（ii）结晶度结论须与 XRD 的相对结晶度一致，若二者矛盾，先检查非晶相
是否真的不吸附 Xe（室温下"可忽略"这一前提）；(iii) Si/Al 与骨架铝的结论须与 $^{27}$Al、
$^{29}$Si MAS NMR 及化学分析对照，尤其注意脱铝产生的非骨架铝物种会使 $\delta_{S}$ 与 Al
含量的关系偏离线性；（iv）阳离子位置与迁移的结论须与离子交换度、ESR（顺磁离子）和
吸附量数据一致；（v）积炭对笼体积的影响见本章后文积炭一节。
\src{S2 §8.4.2，书内 p.547—550 / PDF p.39—42}
\end{crosscheck}
